\documentclass[11pt,leqno]{article}
\usepackage{fontspec}
\usepackage[french,english]{babel}
\usepackage{amsmath}
\usepackage{amssymb}
\usepackage{amsthm}
\usepackage{mathrsfs}
\usepackage{tikz-cd}
\usepackage{enumitem}

\frenchsetup{AutoSpacePunctuation=false}

\newcommand{\printedpage}[1]{\marginpar{\footnotesize p.~#1}}
\newcommand{\msp}[1]{\marginpar{\footnotesize [#1]}}
\newcommand{\DD}{\mathbb{D}}
\newcommand{\cA}{\mathcal{A}}
\newcommand{\Dia}{\mathsf{Dia}}
\newcommand{\Ob}{\mathsf{Ob}}
\newcommand{\Hom}{\mathsf{Hom}}
\newcommand{\hot}{\mathsf{hot}}
\newcommand{\Hot}{\mathsf{Hot}}
\newcommand{\HOT}{\mathsf{HOT}}
\newcommand{\isoto}{\overset{\sim}{\to}}
\newcommand{\isolongto}{\overset{\sim}{\longrightarrow}}
\newcommand{\sHom}{\mathsf{Hom}}
\newcommand{\uHom}{\underline{\mathsf{Hom}}}

\begin{document}

\printedpage{34}
% head: A. Grothendieck : Dérivateurs, Ch. I, Édition des Universités de Montpellier II et Paris VII
% foot: Version 8 juillet 2001
\noindent has contractible fibers.

Another standard situation where base change works (but where the functors involved have no reason to be proper or smooth).

Let
\[
f : I \longrightarrow I' \quad , \qquad g : J \longrightarrow J'
\]
be in $\Dia$, whence the diagram (compare page 23)
\[
\begin{tikzcd}
X' = I \times J \arrow[r, "G'"] \arrow[d, "F'"'] & I \times J' = X \arrow[d, "F"] \\
Y' = I' \times J \arrow[r, "G"] & I' \times J' = Y
\end{tikzcd}
\quad .
\]
This diagram is such that for every left derivator

\noindent\msp{page 60}

\noindent we have base change
\[
G^*F_* \simeq F'_*{G'}^* \quad ,
\]
the reason here being that the induced functors
\[
X'/y' \longrightarrow X/y
\]
are \underline{\underline{co}}strict hot-equivalences, \textit{i.e.} that the categories $K$ (page 49) are contractible. Indeed, if $x_0 = (i_0, j'_0) \in \Ob X$, $y'_0 = (i'_0, j_0) \in \Ob Y'$, $u_0 = (u'_0, v'_0) : F(x_0) \to G(y'_0)$, where $u'_0 : f(i_0) \to i'_0$, $v'_0 : j'_0 \to g(j_0)$ are given, one checks that the category $K$ sought is simply isomorphic to $(i_0, u'_0)\backslash(I/i'_0) \times (j_0, v'_0)\backslash(J/j'_0)$ which is contractible having an initial object \texttt{[rather $(i_0, u'_0)\backslash(I/i'_0) \times (j'_0\backslash J)/(j_0, v'_0)$ which is contractible being the product of a category having an initial object and a category having a final object]}. Whew!

I notice a very big blunder. I thought I remembered that for $f : X \to Y$ which is a strict hot-equivalence, one had for every $F \in \DD(Y)$, $F \isoto f_*f^*F$ (for example for $\DD = \mathsf{Der}(\mathit{Ab})$, or $\DD = \Hot$), \textit{i.e.} $f^*$ fully faithful, whence also that $\Gamma_{\DD}(Y, F) \to \Gamma_{\DD}(X, f^*F)$ is \texttt{[an]} isomorphism. But this is true only if $F$ is a \textit{constant} coefficient on $Y$, \textit{i.e.} of the form $q^*(M)$ ($q : Y \to e$, $M \in \DD(e) = \cA_{\DD}$). This is false, for general $F$, already in the case of the immersion $f : \{\varphi\} \to Y$, where $\varphi$ is an initial object. Then, $f$ is always a strict hot-equivalence, but the homomorphism $\Gamma(Y, F) \to \Gamma(X, f^*F)$ is in general not an isomorphism.

\clearpage
\printedpage{35}
% head: A. Grothendieck : Dérivateurs, Ch. I, Édition des Universités de Montpellier II et Paris VII
% foot: Version 8 juillet 2001
\noindent\msp{page 61}

\subsection*{Notions of $\mathbb{D}$-equivalence and of $\mathbb{D}$-contractibility.}

So we are going to take up again the questions of acyclicity and base change.

\noindent\textbf{Definition 1.} Let $\DD$ be a prederivator, and let
\[
\begin{tikzcd}
X \arrow[rr, "f"] \arrow[dr, "p"'] & & Y \arrow[dl, "q"] \\
& e &
\end{tikzcd}
\]
be in $\Dia$. We say that $f$ is a $\DD$-\textit{equivalence} if the functor
\[
F \longmapsto f^*F \quad , \qquad \DD(Y) \longrightarrow \DD(X)
\]
induces a fully faithful functor on the subcategory of $\DD(Y)$ formed by the \textit{constant} $\DD$-coefficients, \textit{i.e.} if for all $M, N \in \Ob \cA = \DD(e)$ (``absolute coefficients''), we have
\[
\sHom(q^*M, q^*N) \isolongto \sHom(p^*M, p^*N) \quad . \tag*{(1)}
\]
Suppose that $p_*$ and $q_*$ exist. We shall put, for a sheaf $F$ on $X$
\[
H^*_{\DD}(X, F) \overset{\text{def}}{=} p_*(F) \qquad \text{(and likewise } \ H^{\DD}_*(X, F) = p_!(F) \ \text{).}
\]
N.B. And if $F = p^*M$, $M \in \cA$, we write simply $H^*_{\DD}(X, M)$, $H^{\DD}_*(X, M)$.

\noindent The homomorphism (1) is interpreted by adjunction as
\[
\begin{tikzcd}[row sep=small]
\sHom(M, q_*q^*N) \arrow[r] \arrow[d, equal] & \sHom(M, p_*p^*N) \arrow[d, equal] \\
\sHom(M, H^*_{\DD}(Y, q^*N)) \arrow[r] & \sHom(M, H^*_{\DD}(X, p^*N))
\end{tikzcd}
\]
deduced from the homomorphism
\[
q_*q^*N \longrightarrow p_*p^*N \quad , \qquad \textit{i.e.} \quad H^*_{\DD}(Y, N) \longrightarrow H^*_{\DD}(X, N) \quad . \tag*{(2)}
\]
Thus to say that (1) is an isomorphism for all $M, N$ means that (2) is an isomorphism for every $N$. Likewise, if $p_!$, $q_!$ exist this is equivalent to saying that
\[
p_!p^*N \longrightarrow q_!q^*N \quad , \qquad \textit{i.e.} \quad H^{\DD}_*(X, N) \longrightarrow H^{\DD}_*(Y, N) \tag*{$(2')$}
\]
is an isomorphism in $\cA$ for every $N \in \cA$.

\clearpage
\printedpage{36}
% head: A. Grothendieck : Dérivateurs, Ch. I, Édition des Universités de Montpellier II et Paris VII
% foot: Version 8 juillet 2001
\noindent\msp{page 62}

\noindent In summary:

\noindent\textbf{Proposition 1.} \textit{If $p_*$ and $q_*$ exist, i.e. if $H^*_{\DD}(X, \bullet)$ and $H^*_{\DD}(Y, \bullet)$ exist (it even suffices that this exist for constant coefficients), then $f$ is a $\DD$-equivalence if and only if $f$ induces an isomorphism for the $\DD$-cohomologies with constant coefficients of $X$ and $Y$. If $p_!$ and $q_!$ exist, i.e. if $H^{\DD}_*(X, \bullet)$ and $H^{\DD}_*(Y, \bullet)$ exist, $f$ is a $\DD$-equivalence if and only if $f$ induces an isomorphism for the $\DD$-homologies with constant coefficients of $X$ and $Y$.}

\noindent\textbf{Remark 1.} For $f$ to be \texttt{[a]} $\DD$-equivalence, it \textit{suffices} that for every $F$ \textit{constant} on $Y$ the adjunction morphism $F \to f_*f^*F$ be \texttt{[an]} isomorphism.

\noindent\textbf{Remark 2.} If $f$, $f' : X \to Y$ are in the same connected component of $\uHom(X, Y)$, then $f$ is a $\DD$-equivalence $\Leftrightarrow$ $f'$ a $\DD$-equivalence.

\noindent\textbf{Proposition 2.} \textit{Isomorphisms are $\DD$-equivalences. If $X \overset{f}{\to} Y \overset{g}{\to} Z$, if two among $f$, $g$, $gf$ are $\DD$-equivalences, the third too.} % [sic] parmis -> parmi, spelling

\noindent\textbf{Definition 2.} We say that $X$ in $\Dia$ is $\DD$-\textit{aspheric} (or $\DD$-\textit{contractible}, or $\DD$-\textit{trivial}) if $X \overset{p}{\to} e$ is a $\DD$-equivalence, \textit{i.e.} if $p^*$ is fully faithful.

If $p_*$ exists, this therefore means that for $M \in \Ob(\cA)$
\[
M \overset{\mathrm{adj}}{\isolongto} p_*p^*M \quad , \qquad \textit{i.e.} \quad M \isolongto H^*_{\DD}(X, M) \quad ,
\]
and if $p_!$ exists, this means that for every $M$
\[
p_!p^*M \overset{\mathrm{adj}}{\isolongto} M \quad , \qquad \textit{i.e.} \quad H^{\DD}_*(X, M) \isolongto M \quad .
\]
\textbf{Corollary.} \textit{If $Y$} (resp. $X$) \textit{is contractible, then $f : X \to Y$ is a $\DD$-equivalence if and only if $X$} (resp. $Y$) \textit{is contractible.}

\noindent\textsc{Example}. Let us take $\DD = \HOT$. (We can take $\Dia$ = category of small categories). Then, for $M \in \cA = \Hot$ and $X$ a small category, we have
\[
H^*_{\HOT}(X, M) \simeq \sHom_{\cA}(\hot(X), M) \quad ,
\]
where
\[
\hot : \mathcal{C}\mathit{at} \longrightarrow \cA = \Hot
\]
is the canonical functor which associates to every $X$ its ``homotopy type''. We have
\[
H^{\HOT}_*(X, M) = \hot(X) \times M \quad .
\]

\noindent\msp{page 63}

\noindent We deduce from this that $f : X \to Y$ is a $\HOT$-equivalence (in the sense of definition 1) if and only if it is a $\hot$-equivalence (or ``weak equivalence'') \textit{i.e.} if and only if
\[
\hot(f) : \hot(X) \isolongto \hot(Y)
\]
\printedpage{37}
% head: A. Grothendieck : Dérivateurs, Ch. I, Édition des Universités de Montpellier II et Paris VII
% foot: Version 8 juillet 2001
is an isomorphism.

\noindent\textbf{Proposition 2.} \textit{Let $f : X \to Y$. Equivalent conditions (if $\DD$} \texttt{[is a]} left \textit{derivator}). % Resolved against official supplemental edition p37: a gauche). is upright.

a) \textit{$f$ is a $\DD$-equivalence, and remains so after every base change $Y' \to Y$ which is a local isomorphism (for example for $Y/F \to Y$, $F \in \Ob \widehat{Y}$).}

b) \textit{For every $y \in \Ob Y$, the induced functor $X/y \to Y/y$ is a $\DD$-equivalence.}

c) \textit{For every $y \in \Ob Y$, $X/y$ is $\DD$-contractible.}\footnote[25]{Other equivalent conditions:\par
\indent d) \textit{For every \underline{constant} sheaf $M_Y$ on $Y$, we have $M_Y \isoto f_*f^*M_Y$.}\par
\indent d$'$) \textit{As in} d) \textit{with \underline{locally} constant sheaf.}}

\noindent N.B. We say that $g : Z \to Y$ is a \textit{local isomorphism} if for every $z \in \Ob Z$, put $y = g(z)$, the induced morphism on the categories localized at $z$ and $y$
\[
Z/z \isolongto Y/y
\]
is an isomorphism. For example every open immersion, every induced functor $Y/y \to Y$ (localization functor) are local isomorphisms.

From this a) $\Rightarrow$ b) follows trivially.

b) $\Leftrightarrow$ c). Follows from the corollary on the preceding page, if one knows that $Y/y$ is $\DD$-contractible\footnote[26]{Cf. cor. 1 of \textbf{Der 5} (g)}. This will follow from \textbf{Der 5} below.

\noindent\msp{page 64}

c) $\Rightarrow$ a) ``Passage from local to global''. Since hypothesis c) is stable under base change by open immersion, it suffices to prove that c) implies that $f$ is a $\DD$-equivalence, \textit{i.e.} that for every $F = M_Y$ ($M \in \Ob \cA$), constant coefficient on $Y$, the adjunction homomorphism
\[
M_Y \longrightarrow f_*f^*M_Y = f_*M_X
\]
is \texttt{[an]} isomorphism. It suffices to prove that this is an isomorphism fiber by fiber (\textbf{Der 2}), hence that for every $y \in Y$,
\[
M_y \longrightarrow f_*(M_X)_y
\]
is an isomorphism. By \textbf{Der 4} (g), we have
\[
f_*(M_X)_y \isolongto \Gamma(X/y, M) \quad ,
\]
and the hypothesis that $X/y$ is $\DD$-contractible gives what we want.

\noindent\textbf{Proposition 2$'$.} \textit{Dually, if $\DD$ is a left derivator} \texttt{[rather right]}\textit{, equivalent conditions on $f : X \to Y$:}

a) \textit{$f$ a $\DD$-equivalence,} \texttt{[and]} \textit{remains so under every base change $g : Y' \to Y$ which is a ``colocal isomorphism'' (i.e. for every $y' \in \Ob Y'$, the induced functor on the ``colocalized'' categories $y'\backslash Y' \isolongto y\backslash Y$ (where $y = g(y')$)} \texttt{[is an isomorphism]}\textit{).}

b) \textit{For every $y \in Y$, the induced functor $y\backslash X \to y\backslash Y$ is a $\DD$-equivalence.}

c) \textit{For every $y \in Y$, $y\backslash Y$ is $\DD$-contractible.}

\noindent To prove it, one must use that the $y\backslash Y$ are $\DD$-contractible (cf. cor. 1 of \textbf{Der 5} (d)).

\clearpage
\printedpage{38}
% head: A. Grothendieck : Dérivateurs, Ch. I, Édition des Universités de Montpellier II et Paris VII
% foot: Version 8 juillet 2001
\noindent\msp{page 65}

\noindent\textbf{Definition 3.} Let $\DD$ be a prederivator (perhaps even a left or right derivator). We say that $f$ is a \textit{strict} $\DD$-equivalence if $f$ satisfies condition a) of \texttt{[the]} proposition 2. We say that it is a \textit{costrict} $\DD$-equivalence if $f$ satisfies that of \texttt{[the]} proposition $2'$; we also sometimes say that $f$ is $\DD$-cofinal. (The dual condition, \textit{i.e.} strict $\DD$-equivalence, would be called co-$\DD$-cofinal!)\footnote[27]{Better terminology: $f$ is $\DD$-\textit{aspheric} (strict $\DD$-equivalence), $\DD$-\textit{coaspheric} (co-strict $\DD$-equivalence).} % [sic] s'appelerait -> s'appellerait, spelling

\noindent\textbf{Proposition 3.} \textit{Let $X \overset{f}{\to} Y \overset{g}{\to} Z$. If $f, g$ are strict $\DD$-equivalences, then $gf$ too. If $gf$ and $f$ are strict $\DD$-equivalences, then $g$ too. If $\DD$} \texttt{[is a]} \textit{left or right derivator, every equivalence} \texttt{[of categories]} \textit{is a strict $\DD$-equivalence.} (For it has a right adjoint.)

We use the fact that the notion of local isomorphism is stable under base change.

\noindent\msp{page 66}

\noindent\textbf{Proposition 4.} \textit{Let $\DD$ be a right derivator, $f : X \to Y$ a functor in} $\Dia$\textit{, such that for $y \in Y$, $y\backslash X$ is $\DD$-contractible (hence $f$ a costrict $\DD$-equivalence, subject to $\DD$ being a right derivator \ldots\ ). Then for every $F$ in $\DD(Y)$ (not necessarily with constant coefficients) we have} % [sic] constant -> constants, agreement
\[
H^*_{\DD}(Y, F) \isolongto H^*_{\DD}(X, f^*(F))
\]
\textit{is an isomorphism. Conversely, if the condition is true for every $F$, then the $y\backslash X$ are $\DD$-contractible.}

\noindent\textsc{Proof}. It amounts to the same thing to say that for every $M \in \Ob \cA = \DD(e)$, we have
\[
\sHom(M, H^*_{\DD}(Y, F)) \overset{\mathrm{iso}}{\longrightarrow} \sHom(M, H^*_{\DD}(X, f^*(F))) \ .
\]
But by adjunction, this arrow is interpreted as
\[
\sHom(M_Y, F) \longrightarrow \sHom(M_X, f^*F) \tag*{$(*)$}
\]
On the other hand, if $f_!$ exists, we have
\[
\sHom(M_X, f^*F) = \sHom(f_!(M_X), F) \ ,
\]
and the homomorphism $(*)$ is deduced from the adjunction homomorphism
\[
f_!(M_X) = f_!(f^*(M_Y)) \longrightarrow M_Y \ , \tag*{$(**)$}
\]
by taking the $\sHom$ into $F$. Thus the condition is equivalent to saying that this homomorphism is \texttt{[an]} isomorphism, whatever the constant coefficient $M_Y$ on $Y$. To see whether this is an isomorphism, we look fiber by fiber (\textbf{Der 2}), now the fibers are interpreted as $H^{\DD}_*(y\backslash X, M_X)$ (by \textbf{Der 4} (d)) and as $M$.

\end{document}
